\documentclass[text06.tex]{subfiles}
\begin{document}
\newpage
\section{HSPからJuliusを使う}
認識した声を使って、何か動かしたいことがあります。例えば、認識した言葉に\ruby{応}{おう}じて部屋の電気をつけたり、タイマーをセットしたりできます。そんなときには、自分のプログラムとJuliusをつなげて使います。
\subsection{JuliusとHSPの\ruby{連携}{れん|けい}のしくみ}
JuliusとHSPは、別々のプログラムとして動いています。図\ref{JuliusとHSPの連携}のように、Juliusが聞き取った言葉（認識結果）がHSPへ渡されます。この教材では、\code{julius.as} が2つの間で認識結果を受け渡してくれます。今回は、通信の細かい仕組みを覚える必要はありません。
\begin{figure}[H]
\begin{center}
    \includesvg[width=\linewidth]{../../asset/image/JuliuswithHSP.svg}
    \caption{JuliusとHSPの連携}
    \label{JuliusとHSPの連携}
\end{center}
\end{figure}

\subsection{認識した単語を取得・表示する}
\mbox{julius.hsp} を見ながら、まず、特に大事な次の4つをおさえましょう。

\begin{itemize}
\item \code{init\_julius}：音声認識を使う準備をする
\item \code{get\_word\_list}：認識結果を受け取る
\item \code{words(cnt)}：認識された言葉
\item \code{cm(cnt)}：Juliusが、その認識結果にどのくらい自信を持っているかを表す目安
\end{itemize}

まず、これらを使えるようにするために、2行目のように\\
\code{\#include “julius.as”}\\
を書きます。続いて、\code{init\_julius} で音声認識を使う準備をします。引数には辞書ファイル（今回は \mbox{\code{kudamono.dic}}）を指定してください。\code{init\_julius} は、\code{get\_word\_list} を使う前に1回だけ実行します。これによってJuliusが起動します。

プログラムを読むときには、\code{is\_recieved} も使います。Juliusから認識結果が届いたかを調べ、届いていれば0が返ってきます。そのときだけ \code{get\_word\_list} で結果を受け取ります。\\
\code{get\_word\_list words, cm}\\
とすると、認識された言葉が \code{words(cnt)} に、自信の目安が \code{cm(cnt)} に入ります。\code{cnt} は、今見ているのが何番目の言葉かを表します。

\code{cm} は、Juliusが「この認識結果にどのくらい自信を持っているか」を表す目安です。1に近いほど自信が高い結果ですが、正解する\ruby{確率}{かく|りつ}そのものではありません。たとえば、0.8だからといって、80\%の確率で正しい、という意味ではありません。

\begin{lstlisting}[caption=julius.hsp,label=julius.hsp]
#include "hsp3dish.as"
#include "julius.as"

redraw 0
pos 0, 0
mes "Loading Julius..."
redraw 1

init_julius "kudamono.dic"

sockidx = 0
domain = "127.0.0.1"
PORT = 10500

sockopen sockidx, domain, PORT

if stat != 0 {
  redraw 0
  mes "Error in opening socket"
  redraw 1
  wait 30
  stop
}

<#blue#;認識結果を入れておく入れ物。今回は変更しません。#>
sdim words, 4096, 0
ddim cm, 0

repeat -1
  if is_recieved(sockidx) = 0 {
    get_word_list words, cm
    len =  "L1 " + length(words)
    repeat length(words)
      if cm(cnt) > 0.5 {
          ans =  "" + cnt + ": " + words(cnt) + " CM = " + cm(cnt)
      }
    loop
  }
  redraw 0
  pos 0,0
  mes len
  mes ans
  redraw 1
  await 30
loop
\end{lstlisting}

途中でエラーと出たときは、画面にメッセージを出して終了します。

\begin{itembox}[c]{もっと知りたい人へ}
JuliusとHSPは、「ソケット」という通信のしくみでつながっています。\mbox{\code{julius.hsp}} の \code{sockopen} は、自分の\mbox{Raspberry Pi}（127.0.0.1）の10500番ポートでJuliusと通信するための準備です。\code{is\_recieved(sockidx)} の \code{sockidx} は、その通信の番号です。今回は、これらの数字を変えなくても動きます。
\end{itembox}

\begin{tcolorbox}[title=\useOmetoi]
\begin{enumerate}
\addex{ターミナルを開いて、 \textasciitilde /07/に移動しましょう。}
\addex{HSPで\mbox{julius.hsp}を実行しましょう。\\果物の名前を３つ話しかけて、認識されるか確かめてみましょう。}
\end{enumerate}
\end{tcolorbox}
\end{document}
